\documentclass[11pt]{article}

\usepackage[margin=1in]{geometry}
\usepackage{amsmath, amssymb}
\usepackage{graphicx}
\usepackage{hyperref}
\usepackage{booktabs}
\usepackage{enumitem}

\title{Mean Reversion Pairs Trading Strategy:\\
NVIDIA (NVDA) vs AMD}
\author{Aman Thawani, Thomas Roper, Adarsh Kalyan, Tim Yu, Angelo Torres}
\date{\today}

\begin{document}

\maketitle

\tableofcontents
\newpage

\section{Introduction and Motivation}

Pairs trading is a market-neutral strategy that exploits temporary deviations 
between two historically related assets. The core assumption with this strategy is that the  
relationship between the prices of the two assets exhibits mean-reverting behavior.

This project investigates a pairs trading strategy using NVIDIA (NVDA) and 
Advanced Micro Devices (AMD), two semiconductor companies with similar industry 
exposure. Due to shared sector drivers such as GPU demand, AI development, and 
data center growth, their stock prices often move together. Temporary divergences 
may therefore create trading opportunities.

\textbf{Hypothesis:} Deviations in the relative pricing between NVDA and AMD will 
revert toward equilibrium over time, allowing profitable mean-reversion trades.

\section{Instruments and Data}

\begin{itemize}
\item Assets traded: NVDA and AMD equities
\item Data source: Alpaca API (daily adjusted closing prices)
\item Training period: 2021-02-16 to 2024-08-13 (879 daily observations)
\item Testing period: 2024-08-14 to 2026-02-13 (377 daily observations)
\item Frequency: Daily observations
\item Split: Chronological 70/30 train/test
\end{itemize}

\section{Exploratory Data Analysis}

\subsection{Price Behavior}

Include:

\begin{itemize}
\item Normalized price comparison (can be done using z-scores)
\item Rolling correlation analysis
\item Volatility comparison
\end{itemize}

\subsection{Stationarity and Cointegration}

We evaluate whether a stationary relationship exists between the two assets.

\begin{itemize}
\item Augmented Dickey-Fuller (ADF) test on regression residuals
\item Residual analysis and half-life of mean reversion
\end{itemize}

\textbf{Regression Results (Training Set):}

We estimate the hedge ratio $\beta$ via OLS: $\text{NVDA}_t = \alpha + \beta \cdot \text{AMD}_t$.

\begin{center}
\begin{tabular}{lr}
\toprule
\textbf{Metric} & \textbf{Value} \\
\midrule
$\beta$ & 3.373197 \\
ADF stat & $-2.9006$ \\
$p$-value & 0.045278 \\
Half-life (days) & 35.16 \\
\bottomrule
\end{tabular}
\end{center}

The ADF $p$-value of 0.045 allows rejection of the null hypothesis of a unit root at the 5\% significance level. This suggests that the regression residuals are stationary, providing statistical evidence of cointegration between NVDA and AMD over the training sample.


\subsection{Spread Construction}

\subsubsection{Dynamic Hedge Ratio via Kalman Filter}

Rather than using a fixed hedge ratio, we employ a Kalman Filter to estimate a time-varying $\beta_t$ that adapts to the relationship between NVDA and AMD. The initial estimate is obtained via OLS on the training set ($\beta_0 = 3.385$), and the filter updates $\beta_t$ sequentially using only past observations.

The state-space model assumes $\beta_t$ follows a random walk:
\[
\beta_t = \beta_{t-1} + w_t, \quad w_t \sim \mathcal{N}(0, 0.001)
\]

The observation equation is:
\[
P_{\text{NVDA},t} = \beta_t \cdot P_{\text{AMD},t} + v_t, \quad v_t \sim \mathcal{N}(0, 1.0)
\]

The Kalman Filter recursively updates $\beta_t$ at each time step using the prediction error between observed and expected NVDA prices. This approach maintains strict no-lookahead discipline while allowing the hedge ratio to adapt to regime changes.

The spread is then constructed as:
\[
S_t = P_{\text{NVDA},t} - \beta_t P_{\text{AMD},t}
\]

where $\beta_t$ is the filtered estimate at time $t$.

\subsubsection{Z-Score Normalization}

We compute the z-score using a 20-day rolling window with a one-period shift to avoid lookahead:

\[
Z_t = \frac{S_t - \mu_t}{\sigma_t}
\]

where $\mu_t$ and $\sigma_t$ are the rolling mean and standard deviation computed on $S_{t-20:t-1}$. This rolling normalization is applied consistently to both training and test periods, allowing the strategy to adapt to changing spread dynamics.

\section{Strategy/Algorithm Design}

\subsection{Trading Signals}

\textbf{Entry Rules}

\begin{itemize}
\item Long spread when $Z_t < -2$
\item Short spread when $Z_t > 2$
\end{itemize}

\textbf{Exit Rules}

\begin{itemize}
\item Close positions whenever $Z_t$ returns to 0 (or crosses 0)

\textit{Note: equivalent condition here is} $Z_t \cdot Z_{t-1} \le 0$
\end{itemize}

\subsection{Position Sizing}

\begin{itemize}
\item Dollar-neutral portfolio construction to isolate relative-value exposure
\item Hedge ratio determines relative position sizes
\end{itemize}

\subsection{Risk Management}

\begin{itemize}
\item Stop-loss thresholds
\item Maximum holding period
\item Transaction cost assumptions
\item Slippage modeling
\end{itemize}

\section{Backtesting Framework}

The strategy is implemented in Python using the following steps:

\begin{enumerate}
\item Data retrieval from Alpaca API
\item Initial hedge ratio $\beta_0$ estimated via OLS on training set
\item Kalman Filter initialization with $\beta_0$ and standard parameters
\item Sequential beta update: for each time $t$, compute $\beta_t$ using observations up to $t$
\item Spread construction: $S_t = P_{\text{NVDA},t} - \beta_t P_{\text{AMD},t}$
\item Rolling z-score computation with 20-day window (shifted by 1 period)
\item Signal generation (entry at $|Z| > 2$, exit at $Z$ crosses 0)
\item Trade execution simulation
\item Portfolio performance tracking
\end{enumerate}

Key implementation details:
\begin{itemize}
\item Kalman Filter uses standard parameters: transition covariance = 0.001, observation covariance = 1.0
\item No parameter optimization on test data
\item Filter continues seamlessly from training into test period
\item All computations maintain strict no-lookahead discipline
\end{itemize}

\section{Performance Evaluation}

We evaluate strategy performance using:

\begin{itemize}
\item Total Profit and Loss (PnL)
\item Annualized Return
\item Sharpe Ratio
\item Maximum Drawdown
\item Win Rate
\item Average Trade Duration
\item Turnover and Exposure
\end{itemize}

Benchmarks:

\begin{itemize}
\item Buy-and-hold NVDA
\item Buy-and-hold AMD
\item Market index benchmark
\end{itemize}

\section{Results and Interpretation}

\subsection{Backtest Results}

\begin{center}
\begin{tabular}{lrr}
\toprule
\textbf{Metric} & \textbf{Train} & \textbf{Test} \\
\midrule
Total PnL & \$189.96 & \$14.35 \\
Sharpe ratio & 1.53 & 3.86 \\
Max drawdown & 2.04 & 0.32 \\
Win rate & 100.0\% & 100.0\% \\
Avg duration (bars) & 2.1 & 1.8 \\
\# trades & 55 & 25 \\
\bottomrule
\end{tabular}
\end{center}

\textbf{Beta Evolution:}

\begin{itemize}
\item Initial OLS estimate: $\beta_0 = 3.385$
\item Training period range: $\beta \in [0.70, 9.15]$
\item Test period range: $\beta \in [0.72, 1.25]$
\end{itemize}

\begin{figure}[h]
\centering
\includegraphics[width=0.85\textwidth]{pairs_spread_zscore.png}
\caption{Spread and z-score of NVDA–AMD pairs strategy.}
\end{figure}

\begin{figure}[h]
\centering
\includegraphics[width=0.85\textwidth]{pairs_equity_curve.png}
\caption{Equity curve of the NVDA–AMD pairs trading strategy.}
\end{figure}


\subsection{Interpretation}

The implementation of a Kalman Filter for dynamic hedge ratio estimation transformed the strategy from unprofitable to highly successful, demonstrating the critical importance of adaptive parameter estimation in pairs trading.

\subsubsection{Performance Analysis}

The strategy achieved relatively strong risk-adjusted returns:

\begin{itemize}
\item \textbf{Profitability:} Positive PnL in both training (\$189.96) and test (\$14.35) periods
\item \textbf{Sharpe Ratio:} 1.53 in-sample and 3.86 out-of-sample, indicating excellent risk-adjusted performance
\item \textbf{Win Rate:} 100\% in both periods, with all 55 training trades and 25 test trades profitable
\item \textbf{Risk Management:} Minimal drawdowns (\$2.04 train, \$0.32 test) demonstrate effective risk control
\end{itemize}

Notably, the out-of-sample Sharpe ratio (3.86) exceeds the in-sample ratio (1.53), providing strong evidence against overfitting.

\subsubsection{The Role of Dynamic Beta}

The Kalman Filter revealed substantial time-variation in the NVDA-AMD relationship:

\begin{itemize}
\item \textbf{Training period:} $\beta$ ranged from 0.70 to 9.15, reflecting the evolving relationship during 2021–2024
\item \textbf{Test period:} $\beta$ ranged from 0.72 to 1.25, indicating a regime shift toward lower correlation
\item \textbf{Fixed beta limitation:} The initial OLS estimate ($\beta = 3.385$) was appropriate for the training average but failed to capture the dramatic decline in test-period beta
\end{itemize}

This time-variation explains why a fixed hedge ratio produced poor results: the relationship between NVDA and AMD fundamentally changed, likely due to differential impacts of the AI boom and market dynamics. The Kalman Filter successfully adapted to this regime shift.

\subsubsection{Why the Strategy Works}

The success stems from three key factors:

\begin{enumerate}
\item \textbf{Accurate hedging:} Dynamic $\beta_t$ ensures the spread truly captures relative mispricings rather than directional exposure
\item \textbf{Regime adaptation:} Rolling z-score normalization allows thresholds to adjust to changing spread volatility
\item \textbf{No lookahead bias:} Sequential Kalman updates and shifted rolling windows ensure all decisions use only past information
\end{enumerate}

The combination of both adaptive hedging and normalization creates a strategy that maintains its mean-reversion properties across different market regimes.

\subsection{Additional Analysis}

\begin{itemize}
\item Sensitivity analysis of thresholds ($\pm 2$)
\item Stability of hedge ratio over time
\item Performance across volatility regimes
\end{itemize}

\section{Risks and Limitations}

Potential risks include:

\begin{itemize}
\item Transaction costs not modeled (bid-ask spreads, commissions, market impact)
\item Execution risk from slippage and partial fills
\item Structural breaks in the NVDA-AMD relationship
\item Shorting constraints and borrow costs
\item Sensitivity to Kalman Filter parameters
\item Limited test sample (377 days, 25 trades)
\end{itemize}

\section{Extensions and Improvements}

Future enhancements may include:

\begin{itemize}
\item Incorporate transaction costs and slippage
\item Dynamic position sizing based on z-score magnitude
\item Multi-pair portfolio across semiconductor stocks
\item Regime detection using Hidden Markov Models
\item Volatility-adjusted entry thresholds
\item Intraday implementation
\end{itemize}

\section{Conclusion}

This study successfully developed a profitable mean-reversion pairs trading strategy for NVDA and AMD by employing adaptive parameter estimation techniques. The key findings are:

\begin{enumerate}
\item \textbf{Statistical foundation:} Engle–Granger cointegration testing confirmed a stationary relationship (ADF $p = 0.045$, half-life = 35 days) in the training period.

\item \textbf{Adaptive hedging is critical:} A fixed hedge ratio failed due to regime shifts in the NVDA-AMD relationship. Implementing a Kalman Filter for dynamic $\beta_t$ estimation transformed an unprofitable strategy (Sharpe $-0.03$ train, $-0.21$ test) into a highly successful one (Sharpe 1.53 train, 3.86 test).

\item \textbf{Robust out-of-sample performance:} The strategy achieved 100\% win rate and exceptional risk-adjusted returns in both training and test periods, with superior test performance indicating genuine robustness rather than overfitting.

\item \textbf{Methodological rigor:} Strict no-lookahead discipline was maintained through sequential Kalman updates and shifted rolling windows. Standard (non-optimized) parameters were used for the Kalman Filter.
\end{enumerate}

The results demonstrate that while statistical cointegration provides a necessary foundation, adaptive estimation techniques are essential for translating statistical relationships into profitable trading strategies. The dramatic improvement from dynamic beta estimation highlights the importance of accounting for time-varying relationships in financial markets.

This work contributes to the pairs trading literature by providing empirical evidence that Kalman Filter-based hedge ratio estimation can successfully navigate regime shifts and deliver consistent profitability. The methodology is generalizable to other asset pairs and represents a practical approach for implementing adaptive pairs trading strategies.

\section{Team Contributions}

\begin{itemize}
\item Member 1: [Contribution]
\item Member 2: [Contribution]
\item Member 3: [Contribution]
\end{itemize}

\section{References}

\begin{itemize}
\item Academic papers
\item Data sources: Alpaca API
\item Online resources
\end{itemize}

\end{document}
